% Generated draft: requires mathematical and source review.
\documentclass[11pt]{article}
\usepackage[margin=25mm]{geometry}
\usepackage{fontspec}
\setmainfont{Latin Modern Roman}
\usepackage{amsmath,amssymb,mathtools,unicode-math}
\setmathfont{Latin Modern Math}
\usepackage{enumitem,needspace,xurl,hyperref}
\setlength{\parindent}{0pt}
\setlength{\parskip}{5pt}
\setlength{\emergencystretch}{3em}
\newcommand{\N}{\mathbb{N}}
\newcommand{\Z}{\mathbb{Z}}
\newcommand{\Q}{\mathbb{Q}}
\newcommand{\R}{\mathbb{R}}
\newcommand{\C}{\mathbb{C}}
\newcommand{\F}{\mathbb{F}}
\newcommand{\A}{\mathbb{A}}
\newcommand{\PP}{\mathbb{P}}
\newcommand{\E}{\mathbb{E}}
\newcommand{\eps}{\varepsilon}
\newcommand{\dd}{\,\mathrm{d}}
\begin{document}
\section*{UnsolvedMath ID: 30000669}
\textbf{Source problem code:} catalogue entry
\subsection{Definitions and mathematical statement}
Let a connected reductive complex algebraic group $G$ act on $\C^n$, and let its two-dimensional affine categorical quotient be
\[
 \C^n//G=\operatorname{Spec}\C[\C^n]^G
                  \simeq\C^2/\Gamma
\]
for a finite algebraic action of $\Gamma$. The source concludes that $\Gamma$ should be cyclic, meaning generated by one element. Connectedness is already in its first hypothesis and is repeated in its concluding clause; it is retained.

\textbf{The unrestricted presentation of $\Gamma$ is insufficient.} Take $G=\mathbb G_m$ acting effectively on $\C^3$ by $a\cdot(x,y,z)=(x,y,az)$. Then $\C[\C^3]^G=\C[x,y]$ and the quotient is $\C^2$. On the other hand the noncyclic group $S_3$ acts faithfully on
\[
 V=\{(u_1,u_2,u_3)\in\C^3:u_1+u_2+u_3=0\}\simeq\C^2
\]
by permutation. Its invariant ring is $\C[e_2,e_3]$, where $e_2,e_3$ are the elementary symmetric polynomials restricted to $V$, so $V/S_3\simeq\C^2$. This meets the extracted quotient hypotheses with a noncyclic finite group, even with faithful actions.

Thus a finite quotient presentation is not unique in a way that permits the literal cyclicity conclusion. One possible intended statement could concern a distinguished small/no-pseudoreflection quotient group, or existence of \emph{some cyclic} quotient presentation rather than cyclicity of every given presentation. These are different assertions and neither missing restriction is supplied. The example uses complex reflections; no condition excluding them or requiring an isolated nontrivial quotient singularity occurs in the record.

The exact original claim is preserved for source review instead of silently replacing $\Gamma$ by a preferred local group. The literature triage checked 21 August 2026 reports no authoritative complete resolution, but does not resolve this presentation ambiguity or elementary counterexample. That open label is not treated as certification of the literal assertion or as a status assessment of any repaired conjecture.
\subsection{Short English statement}
The source predicts a cyclic finite group for a two-dimensional quotient of affine space by a connected reductive group. As written, arbitrary finite presentations invalidate that conclusion: a faithful $S_3$ reflection quotient is already $\C^2$. A canonical/small-group convention or an existence-of-cyclic-presentation target is missing.
\subsection{Sources}
\begin{itemize}
\item Local UnsolvedMath registry, permanent ID 30000669, exact source code catalogue entry; \emph{Cyclicity of Surface Quotients by Connected Groups}. Proposed year 2007, status open; corrected\_verified extraction preserves the full claim but supplies no faithfulness, small-group or pseudoreflection convention for $\Gamma$.
\item \emph{Affine Algebraic Geometry (2007)}, Oberwolfach Reports, DOI 10.4171/owr/2007/01, pp. 5--82, \url{https://doi.org/10.4171/owr/2007/01}.
\item Supplied literature triage checked 2026-08-21 cites the same report, \url{https://doi.org/10.4171/owr/2007/01}, and reports no authoritative complete resolution found. The dated OWR-derived curation requests independent verification of time-sensitive status and cited claims; no canonical quotient-group definition is supplied.
\end{itemize}
\end{document}
